\begin{abstract}
Este artículo presenta un análisis comparativo exhaustivo entre las arquitecturas de combustores anulares y can-anulares en motores de aviación comercial. Se evalúan métricas clave de desempeño como eficiencia de combustión, distribución de temperatura, pérdida de presión y estabilidad de operación bajo condiciones representativas de vuelo. Desde la perspectiva de sostenibilidad, se cuantifican emisiones de NOx, CO, HC y partículas, junto con el potencial de integración con combustibles sostenibles de aviación (SAF) y hidrógeno. Utilizando datos de simulaciones CFD, pruebas en banco y revisiones de literatura reciente (2022-2026), se identifican ventajas de los combustores anulares en uniformidad térmica y reducción de emisiones, frente a la mayor mantenibilidad de los can-anulares. Los resultados indican reducciones potenciales de hasta 40% en NOx y mejoras en eficiencia del 5-10% con diseños anulares optimizados. Se discuten implicaciones para el diseño de próxima generación de motores turbofan de alto bypass ratio y se proponen recomendaciones para futuras investigaciones en combustión lean y materiales avanzados.
\end{abstract}